\documentclass[11pt]{letter}

\usepackage[margin=0.75in]{geometry}
\usepackage{setspace}
\usepackage{hyperref}
\usepackage{xcolor}
\usepackage{enumitem}
\usepackage{parskip}

\hypersetup{
    colorlinks=true,
    urlcolor=black,
    linkcolor=black
}

\setlength{\parindent}{0pt}
\setstretch{1.05}
\date{}
\begin{document}

\begin{letter}{}

\opening{Dear Hiring Committee,}

I saw the job post on \textbf{the Pariveda careers website} and thought I'd be a good fit for it.

I am a Master's student in Computer Science at the University of Cincinnati, graduating December 2026, looking for a role as an \textbf{Entry-Level Software Engineer} at \textbf{Pariveda}. With hands on experience building backend automation, full stack applications, cloud infrastructure, and Generative AI pipelines, I am excited about the opportunity to contribute to \textbf{Pariveda} and support client teams solving real business challenges through technology.

\textbf{Here are reasons that make me a great fit for the role:}

\begin{itemize}[leftmargin=*, itemsep=0.45em]
    \item \textbf{Software development experience across the stack:} At Siemens Industry Software, I automated Teamcenter X infrastructure workflows with Python (Boto3) on AWS, replacing legacy Ruby scripts and cutting deployment time by 40 percent. I also executed 44 production environment upgrades in 4 months, diagnosing and resolving live infrastructure issues across Linux, EC2, databases, and application services, which gave me practical experience estimating, developing, testing, and delivering under real production constraints.

    \item \textbf{Cloud and DevOps depth across AWS, CI/CD, Docker, and Terraform:} I built CI/CD pipelines and Kubernetes deployments using GitLab CI/CD, ArgoCD, and Helm, migrating 40 plus production environments and designing a Terraform and Ansible based Blue Green deployment framework that reduced upgrade downtime by 50 percent, directly matching the cloud and DevOps practices Pariveda consultants work with daily.

    \item \textbf{Generative AI and LLM experience:} As a Graduate Researcher, I designed a two stage memory augmented LLM agent pipeline using Google Gemini and Vertex AI for Android bug reproduction, achieving a 25.7x latency reduction over a GPT-4o based baseline while preserving accuracy. I also built an evidence grounded RAG pipeline over Gemini 2.5-flash with a 17 tool MCP server for a healthcare eligibility project, giving me direct experience with RAG, prompt engineering, and agent frameworks.

    \item \textbf{Strong communication and collaborative problem solving:} As Quality Coach and SonarQube Specialist at Siemens, I increased automated code coverage from 35 to 78 percent and mentored 5 software engineers through structured onboarding, work that required explaining technical tradeoffs clearly to both engineers and stakeholders, a skill I look forward to applying when presenting solutions to clients.
\end{itemize}

I am available for full time work after my graduation in December 2026. As a Master's student at the University of Cincinnati and someone who values professionalism and growth, I am enthusiastic about contributing to Pariveda's mission of helping clients get better at what they do, from idea to impact.

Please find my CV attached below. I look forward to hearing from you. I would be excited to grow as a technologist and future leader within a team that values continuous learning and using technology to help others succeed.

{\setlength{\parskip}{0pt}
Regards,\\
Tanmay Bhuskute,\\
bhuskutv@mail.uc.edu $\mid$ tanmay.v.bhuskute@gmail.com $\mid$ +1 7328583027,\\
3424 Brookline Ave, Flat 19, Cincinnati, OH 45220.
}

\end{letter}

\end{document}
